% Generated from project outputs. Do not edit by hand.
\begin{table}[!htbp]
\centering
\begin{threeparttable}
\caption{Including remittances by subsample: 2016-2024}
\label{tab:descriptive-2016-2024-subsamples-y4}
\small
\setlength{\tabcolsep}{4pt}
\renewcommand{\arraystretch}{1.15}
\begin{tabularx}{\linewidth}{@{}Xrrrrr@{}}
\toprule
Sample & N & Mean t & SD t & Mean t+1 & SD t+1 \\
\midrule
General & 552 & 2,970.8 & 1,798.4 & 3,175.9 & 1,809.5 \\
Men & 552 & 3,305.1 & 2,098.5 & 3,537.8 & 2,083.8 \\
Women & 552 & 2,420.6 & 1,581.9 & 2,615.3 & 1,664.3 \\
Urban & 552 & 3,206.1 & 1,846.0 & 3,390.6 & 1,830.1 \\
Rural & 547 & 2,466.2 & 1,663.5 & 2,727.3 & 1,881.2 \\
Indigenous & 547 & 2,439.3 & 1,706.2 & 2,710.1 & 1,783.4 \\
Non-Indigenous & 460 & 3,135.4 & 1,875.2 & 3,273.9 & 1,830.1 \\
Bottom 99 percent & 552 & 2,750.2 & 1,424.1 & 2,950.9 & 1,448.9 \\
Bottom 90 percent & 552 & 2,285.0 & 1,002.1 & 2,444.7 & 1,001.6 \\
Bottom 50 percent & 548 & 1,329.9 & 565.4 & 1,436.2 & 562.6 \\
\bottomrule
\end{tabularx}
\begin{tablenotes}[flushleft]
\footnotesize
\item Monthly income in quetzales. Unweighted descriptive statistics of survey-weighted cohort means. N counts cohort-transition rows, not individuals or independent cohorts. SD is dispersion across cohort means, not the standard error of a mean. Six annual transitions are pooled; the 2019-2021 gap is excluded. Both incomes must be observed. Subsamples overlap; bottom-income groups are cumulative.
\end{tablenotes}
\end{threeparttable}
\end{table}
